% Ad Atticum 13.16 (ad-atticum-13-16)
% Source: https://raw.githubusercontent.com/PerseusDL/canonical-latinLit/master/data/phi0474/phi057/phi0474.phi057.perseus-lat1.xml
% Fetched by scripts/fetch_latin.py

% Dateline (Perseus): Scr. in Arpinati iv K. Quint. a. 709 (45) .

\ciceroLetterOpener{CICERO ATTICO salutem}

\ciceroSection{1} nos cum flumina et solitudinem sequeremur quo facilius sustentare nos possemus, pedem e villa adhuc egressi non sumus; ita magnos et adsiduos imbris habebamus. illam Ἀκαδημικὴν σύνταξιν totam ad Varronem traduximus. primo fuit Catuli, Luculli, Hortensi; deinde quia παρὰ τὸ πρέπον videbatur, quod erat hominibus nota non illa quidem ἀπαιδευσία sed in iis rebus ἀτριψία , simul ac veni ad villam, eosdem illos sermones ad Catonem Brutumque transtuli. ecce tuae litterae de Varrone. nemini visa est aptior Antiochia ratio.

\ciceroSection{2} sed tamen velim scribas ad me, primum placeatne tibi aliquid ad illum, deinde, si placebit, hocne potissimum. quid ? Servilia iamne venit? Brutus ecquid agit et quando? de Caesare quid auditur? ego ad Nonas, quem ad modum dixi. tu cum Pisone, si quid poteris.
